\newpage
\begin{center}
\section{Conclusioni}
\medskip
\end{center}
{\color{gray}\hrule}
\medskip
L'esperienza ha permesso di verificare e in seguito analizzare la conducibilià ottica ed elettrica di un campione di semiconduttore di Seleniuro di Gallio.

Prima di misurare effettivamente le proprietà ottiche ed elettriche del semiconduttore si è
ricavata, per mezzo di uno spettrometro, la curva di calibrazione del monocromatore:
$$\lambda (\#_{ch}) = (  177 \pm 3)\text{nm } + ( 0.388 \pm 0.006 )\frac{\text{nm}}{\#} \cdot \#_{ch}  +  (  -2.2 \pm 0.3 ) \frac{\text{nm}}{\#^2} \cdot \#_{ch}^2$$
$$\#_{ch} (x) = ( -660 \pm 30)\# + ( 3.00 \pm 0.05 )\frac{\#}{\text{u.a.}} \cdot x$$

I dati di trasmittanza ottica e fotocorrente sono in buon accordo con l'andamento caratteristico di un semiconduttore: in accordo con le previsioni teoriche, si è osservato un marcato aumento della fotocorrente in corrispondenza di una diminuzione della trasmittanza ottica per valori di $\lambda$ inferiori a uno stesso $\lambda_g$. 

Dallo studio dell'andamento dei dati, interpolati da due curve sigmoidi, si sono ricavate 
due stime di $\lambda_g$ corrispondenti ai punti di flesso, con le quali si è stimato
un valore dell'\textit{energy gap} pari a $$E_g = (1.962 \pm 0.008)\,\text{eV}$$
che è in buon accordo con l’intervallo atteso sulla base dell'analisi qualitativa fatta in precedenza.


